\PassOptionsToPackage{dvipsnames,svgnames,table}{xcolor}
\documentclass[10pt]{article}
\usepackage[a4paper,margin=22mm]{geometry}
\usepackage[no-math]{fontspec}\setmainfont[SmallCapsFont={Latin Modern Roman Caps}]{Latin Modern Roman}
\usepackage{amsmath,amssymb,amsthm,mathtools,graphicx,xcolor,hyperref,xurl,xspace,ifthen,enumerate,textcomp}
\usepackage[shortlabels]{enumitem}
\setlength{\emergencystretch}{4em}
\newtheorem{theorem}{Theorem}[section]
\newtheorem{theo}[theorem]{Theorem}
\newtheorem{thm}[theorem]{Theorem}
\newtheorem{lemma}[theorem]{Lemma}
\newtheorem{lemm}[theorem]{Lemma}
\newtheorem{prop}[theorem]{Proposition}
\newtheorem{coro}[theorem]{Corollary}
\newtheorem{corollary}[theorem]{Corollary}
\newtheorem{fact}[theorem]{Fact}
\newtheorem{claim}[theorem]{Claim}
\theoremstyle{definition}
\newtheorem{defi}[theorem]{Definition}
\newtheorem{definition}[theorem]{Definition}
\newtheorem{rema}[theorem]{Remark}
\newtheorem{remark}[theorem]{Remark}
\newtheorem{exam}[theorem]{Example}
\newtheorem{example}[theorem]{Example}
\newenvironment{enonce*}[2][]{\par\medskip\noindent\textbf{#2.}\itshape}{\par\medskip}
\providecommand{\eg}{e.g.\ }
\providecommand{\ie}{i.e.\ }
\providecommand{\resp}{resp.\ }
\providecommand{\cf}{cf.\ }
\providecommand{\longthanks}[1]{\par\smallskip\noindent #1\par\smallskip}
\newtheorem{proposition}[theorem]{Proposition}
\newtheorem{conjecture}[theorem]{Conjecture}
\numberwithin{equation}{section}
\newcommand{\R}{{\mathbb R}}
\newcommand{\Z}{{\mathbb Z}}
\newcommand{\N}{{\mathbb N}}
\newcommand{\expref}[2]{\texorpdfstring{\hyperref[#2]{#1~\ref{#2}}}{#1~\ref{#2}}}
\newcommand{\epfmtdoi}[1]{\href{https://doi.org/#1}{doi:#1}}
\newcommand{\epfmt}[2]{#1:\,#2}
\newcommand{\eps}{\varepsilon}
\newcommand{\poly}{\textrm{poly}}
\begin{document}
\begingroup\raggedright
\section*{1598. Conditional Boolean matrix multiplication speedup from an algorithmic triangle-removal lemma}
Nikhil Bansal and Ryan Williams. \emph{Regularity Lemmas and Combinatorial Algorithms}. Theory of Computing volume 8 article 4 (2012), official complete journal PDF acquired 2026-09-30; canonical DOI 10.4086/toc.2012.v008a004 verified against official landing; printed header identity cross-checked and any variant preserved. \url{https://theoryofcomputing.org/articles/v008a004/}. CC BY3.0.
\paragraph{Selected problem and scope (editorial).} Exact original conditional Theorem2.1 only: given the stated algorithmic triangle-removal guarantee T(n,epsilon), f(epsilon), Boolean matrix multiplication has the displayed randomized high-probability runtime with all original word-size, sampling, sparse-product and output terms. Complete Boolean-semiring/word-RAM/pointer-machine model, named sparse-vector preprocessing statements, supplied sparse-product compression explanation and folklore proof, witness sampling, tripartite reduction/three removal-edge classes, full new subset-compressed triangle-reporting construction and all pointer/word implementations and accounting are retained. Named Feder-Motwani clique compression, Blelloch-Vassilevska-Williams preprocessing and elementary binomial/union bounds and word operations are prerequisites; the new reporting graph and all local reduction core are supplied. The algorithmic removal premise is an explicit assumption, not an unconditional new theorem. Weak-regularity speedup, independent-set queries and parameter corollaries are unselected and uncounted.
\paragraph{Difficulty (editorial).} After the named sparse-compression and elementary sampling inputs, the reduction still has to recover all remaining matrix entries from three different removal-edge classes and construct a subset-compressed graph for the otherwise unresolved sparse third class. The complete reporting proof at331–375 proves both an output-sensitive witness budget and a separate word-level implementation, with quantitative dependence on sparsity, witnesses and word size. Those full construction/accounting obligations form one H2 conditional algorithmic reduction above a direct triangle-removal invocation or elementary Four-Russians calculation.
\paragraph{Formatting (editorial).} Exact original human mathematical body and order retained; standalone typography, counter restoration and declared noncore printed locators only. Original comment prose excluded with percent guards retained; source typography and counters restored without upstream class execution. No mathematical repair or reconstruction.
\paragraph{Original source quirks (editorial).} Native321 retains its original sparse-product prose bound without the n factor present in the supplied Theorem3.9 at267 and selected primary bound at137. Native373 retains the original extra n in its word-count prose formula, with the exact preceding loop and full stated reporting bound unchanged. The graph-construction paragraph337 retains k/l versus kappa/ell typography, and native317 retains T(n) while the selected statement uses T(n,epsilon). These original displays and surrounding authored text are preserved without mathematical correction or refereeing.
\endgroup
\clearpage
\setcounter{section}{1}
\section{Our results}

In this paper we present what are arguably the first concrete improvements on combinatorial algorithms for dense BMM since the 70's. Our approach opens a new line of attack on the problem by connecting the complexity of BMM to modern topics in graph theory, such as the Weak Regularity Lemma and the efficiency of Triangle Removal Lemmas.

\subsection{Triangle Removal Lemmas and BMM}

A Triangle Removal Lemma~\cite{RuzsaSz,Green} states that there is a function $f$ satisfying $\lim_{x \rightarrow 0} f(x) = 0$ such that for every graph with at most $\eps n^3$ triangles, we can efficiently find $f(\eps) n^2$ edges that hit all triangles. This lemma is one of the many deep consequences of Szemer\'{e}di's Regularity Lemma~\cite{sz}. We prove that good removal lemmas imply faster Boolean matrix multiplication. Let $w$ be the wordsize (typically $w = \Theta(\log n)$).

\begin{theorem} \label{boolmmtriangles} Suppose there is an $O(T(n,\eps))$ time algorithm that, for every graph $G=(V,E)$ with at most $\eps n^3$ triangles, returns a set $S \subseteq E$ with $|S| \leq f(\eps) n^2$ such that $G'=(V,E\setminus S)$ is triangle-free. Then there is a randomized algorithm for Boolean matrix multiplication that returns the correct answer with high probability and runs in time
%
 \[
O\left(T(n,\eps) + \frac{f(\eps) n^3 \log (1/f(\eps))}{w \log n} + \frac{n^2}{\eps} \cdot \log n + \eps n^3\right)\,.
\]
%
\end{theorem}
\setcounter{section}{2}
\section{Preliminaries}
The \emph{Boolean semiring} is the semiring on $\{0,1\}$ with OR as addition and AND as multiplication. For Boolean matrices $A$ and $B$, $A \vee B$ is the componentwise OR of $A$ and $B$, $A \wedge B$ is the componentwise AND, and $A \star B$ is the (Boolean) matrix product over the Boolean semiring. When it is clear from the context, sometimes we omit the $\star$ and write $A B$ for the product.

Since the running times of our algorithms involve polylogarithmic terms, we must make the computational model precise. Unless otherwise specified, we assume a standard word RAM with wordsize $w$. That is, accessing a memory location takes $O(1)$ time, and we can perform simple operations (such as addition, componentwise AND and XOR, but not multiplication) on $w$-bit numbers in $O(1)$ time.
Typically, speedups in combinatorial algorithms come from either exploiting some combinatorial substructure, by preprocessing and doing table lookups, or by some ``word tricks'' which utilize the bit-level parallelism of the machine model.
In our results, we explicitly state the dependence of the word size, denoted by $w$. The reader may assume $w = \Theta(\log n)$ for convenience. In fact all algorithms in this paper can be implemented on a pointer machine under this constraint.

We now describe some of the tools we need.
\setcounter{subsection}{1}\setcounter{theorem}{6}
\subsection{Preprocessing Boolean matrices for sparse operations}

Our algorithms exploit regularity to reduce dense BMM to a collection of somewhat sparse matrix multiplications. To this end, we need results on preprocessing matrices to speed up computations on sparse inputs. The first deals with multiplication of an arbitrary matrix with a sparse vector, and the second deals with multiplication of a sparse matrix with another (arbitrary) matrix.

\begin{theorem}[Blelloch-Vassilevska-Williams~\cite{ICALP08}] \label{processforsparse}
Let $B$ be a $n \times n$ Boolean matrix and let $w$ be the wordsize. Let $\kappa \geq 1$ and $\ell > \kappa$ be integer parameters. There is a data structure that can be constructed with $O((n^2 \kappa / \ell) \cdot \sum_{b=1}^{\kappa} \binom{\ell}{b})$ preprocessing time, so that for any Boolean vector $v$, the product $B \star v$ can be computed  in 
\[
O\left(n \log n + \frac{n^2}{\ell w} + \frac{n t}{\kappa w}\right)
\]
time, where $t$ is the number of nonzeros in $v$.
\end{theorem}

This result is typically applied as follows. Fix a value of $t$ to be the number of nonzeros we expect in a typical vector $v$. Choose $\ell$ and $\kappa$ such that $n/\ell \approx t/\kappa$, and $\sum_{b=1}^{\kappa} \binom{\ell}{b} = n^\delta$ for some $\delta > 0$. One such choice is $\kappa = \delta \ln(n)/\ln(en/t)$ and $\ell = \kappa \cdot en/t$ in which case we obtain:

\begin{theorem}
Let $B$ be a $n \times n$ Boolean matrix.  There is a data structure that can be constructed with $\tilde{O}(n^{2+\delta})$ preprocessing time, so that for any Boolean vector $v$, the product $B \star v$ can be computed in 
$$
O\left(n \log n + \frac{n t \ln(en/t)}{\delta w \ln n }\right)
$$
time, where $t$ is the number of nonzeros in $v$.
\end{theorem}

We should remark that we do not explicitly apply the above theorem, but the idea (of preprocessing for sparse vectors) is used liberally in this paper.

The following result is useful for multiplying a sparse matrix with another arbitrary matrix.
\begin{theorem}\label{sparsemm} There is an $O(m n \log(n^2/m)/(w\log n))$ time algorithm for computing $A \star B$, for every $n\times n$ $A$ and $B$, where $A$ has $m$ nonzeros and $B$ is arbitrary.
\end{theorem}
This result follows in a straightforward manner by combining the two lemmas below.
The first is a graph compression method due to Feder and Motwani.

\begin{lemma}[From Feder-Motwani~\cite{federmotwani}, Theorem 3.3]\label{sparseramsey} Let $\delta \in (0,1)$ be constant. We can write any $n \times n$ Boolean matrix $A$ with $m$ nonzeros as $A = (C \star D) \vee E$ where $C$, $D$ are $n \times m/n^{1-\delta}$, $m/n^{1-\delta} \times n$, respectively, both with at most $m (\log n^2/m)/(\delta \log n)$ nonzeros, and $E$ is $n \times n$ and has at most $n^{2-\delta}$ nonzeros. Furthermore, finding $C$, $D$, $E$ takes $O(mn^{\delta}\log^2 n)$ time.\end{lemma}

Since the lemma is not stated explicitly in~\cite{federmotwani}, let us sketch the proof for completeness. Using algorithmic Ramsey theoretic arguments (\ie, finding large bipartite cliques in a dense enough graph), Feder and Motwani show that for every bipartite graph $G$ on $2n$ nodes (with $n$ nodes each on left and right) and $m>n^{2-\delta}$ edges, its edge set can be decomposed into $m/n^{1-\delta}$ edge-disjoint bipartite cliques, where the total sum of vertices over all bipartite cliques (a vertex appearing in $K$ cliques is counted $K$ times) is at most $m (\log n^2/m)/(\delta \log n)$. Every $A$ can be written in the form $(C \star D) \vee E$, by having the columns of $C$ (and rows of $D$) correspond to the bipartite cliques. Set $C[i,k]=1$ iff the $i$th node of the LHS of $G$ is in the $k$th bipartite clique, and similarly set $D$ for the nodes on the RHS of $G$. Note that  $E$ is provided just in case $A$ turns out to be sparse.

We also need the following simple folklore result. It is stated in terms of wordsize $w$, but it can easily be implemented on other models such as pointer machines with $w=\log n$.

\begin{lemma}[Folklore]\label{sparseword} There is an $O(m n/w + pq + pn)$ time algorithm for computing $A \star B$, for every $p \times q$ matrix $A$ and $q \times n$  matrix $B$ where $A$ has $m$ nonzeros and $B$ is arbitrary.\end{lemma}

\begin{proof} We assume the nonzeros of $A$ are stored in a list structure; if not we construct this in $O(pq)$ time. Let $B_j$ be the $j$th row of $B$ and $C_i$ be the $i$th row of $C$ in the following. We start with an output matrix $C$ that is initially zero. For each nonzero entry $(i,j)$ of $A$, update $C_i$ to be the OR of $B_j$ and $C_i$. Each update takes only $O(n/w)$ time. It is easy to verify that the resulting $C$ is the matrix product.\end{proof}

\section{Combinatorial Boolean matrix multiplication via triangle removal}

In this section, we prove \expref{Theorem}{boolmmtriangles}.
That is, we show that a more efficient Triangle Removal Lemma implies more efficient Boolean matrix multiplication.
Let $A$ and $B$ be the matrices whose product $D$ we wish to compute. The key idea is to split the task into two cases. First, we use simple random sampling to determine the entries in the product that have many witnesses (where $k$ is a \emph{witness for $(i,j)$} if $A[i,k]=B[k,j]=1$). To compute the entries with few witnesses, we set up a tripartite graph corresponding to the remaining undetermined entries of the matrix product, and argue that it has few triangles. (Each triangle corresponds to a specific witness for a specific entry in $D$ that is still undetermined.) By a Triangle Removal Lemma, a sparse number of edges hit all the triangles in this graph.\footnote{Note that the triangle removal lemma may also return edges that do not lie in any triangle.} Using three carefully designed sparse matrix products (which only require one of the matrices to be sparse), we can recover all those entries $D[i,j]=1$ which have few witnesses.

We now describe our algorithm for BMM.

\smallskip

\noindent{\bf Algorithm:} Let $A$ and $B$ be $n \times n$ matrices. We wish to compute $D = A \star B$, \ie,
\[
D[i,j] = \left(\bigvee_{k=1}^n A[i,k] \wedge B[k,j]\right)\,.
\] 

\smallskip

\noindent
\emph{Random sampling for pairs with many witnesses.}
First, we detect the pairs $(i,j)$ with at least $\eps n$ witnesses.
Construct a $n\times n$ matrix $C$  as follows. Pick a sample $R$ of $(6\log n)/\eps$ elements from $[n]$.
For each $(i,j)$,  $1 \leq i,j \leq n$, check if there is a $k \in R$ that is a witness for $(i,j)$ in the product. If yes, set $C[i,j]=1$, otherwise $C[i,j]=0$. Clearly, this takes at most $O((n^2 \log n) /\eps)$ time. Note that  $C$ is dominated by the desired $D$, in that $C[i,j] \leq D[i,j]$ for all $i,j$.
If $(i,j)$ has at least $\eps n$ witnesses, then some witness lies in $R$ with probability at least $1-1/n^6$. Thus with probability at least $1-1/n^4$, $C[i,j]=D[i,j]=1$ for every $(i,j)$ with at least $\eps n$ witnesses.

\smallskip

\noindent
\emph{Triangle removal for pairs with few witnesses.} It suffices to determine those $(i,j)$ such that $C[i,j] = 0$ and $D[i,j] = 1$. We shall exploit the fact that such pairs do not have many witnesses. Make a tripartite graph $H$ with vertex sets $V_1$, $V_2$, $V_3$, each with $n$ nodes indexed by $1,\ldots,n$. Define edges as follows:
\begin{itemize}
\item Put an edge $(i,k) \in (V_1,V_2)$ if and only if $A[i,k]=1$.
\item Put an edge  $(k,j) \in (V_2,V_3)$  if and only if $B[k,j]=1$.
\item Put an edge $(i,j) \in (V_1,V_3)$ if and only if $C[i,j]=0$.\end{itemize}
That is, edges from $V_1$ to $V_3$ are given by $\overline{C}$, the complement of $C$. Observe that
$(i,k,j) \in (V_1,V_2,V_3)$ is a triangle if and only if $k$ is a witness for $(i,j)$ and $C[i,j] = 0$. Thus our goal is to find the pairs $(i,j)\in (V_1,V_3)$ that are in triangles of $H$.

Since every $(i,j) \in (V_1,V_3)$ has at most $\eps n$ witnesses, there are at most $\eps n^3$ triangles in $H$. Applying the promised Triangle Removal Lemma (in \expref{Theorem}{boolmmtriangles}), we can find in time $O(T(n))$ a set of edges $F$ where $|F| \leq f(\eps) n^2$ and each triangle must use an edge in $F$. Hence it suffices to compute those edges $(i,j) \in (V_1,V_3)$ that participate in a triangle with an edge in $F$. 

Define $A_F[i,j]=1$ if and only if $A[i,j]=1$ and $(i,j) \in F$. Similarly define $B_F$ and $\overline{C}_F$. Every triangle of $H$ passes through at least one edge from one of these three matrices. Let $T_A$ (resp. $T_B$ and $T_{\overline{C}}$) denote the set of triangles with an edge in ${A}_F$ (resp. $B_F$ and ${\overline{C}}_F$).  Note that we do not know these triangles.

We can determine the edges $(i,j) \in (V_1,V_3)$ that are in some triangle in $T_A$ or $T_B$ directly by computing $C_1 = {A}_F \star B$ and $C_2 = A \star {B}_F$, respectively. As $A_F$ and $B_F$ are sparse, by \expref{Theorem}{sparsemm}, these products can be computed in $O(|F|\log(n^2/|F|)/(w\log n))$ time. The $1$-entries of $\overline{C} \wedge C_1$ (resp. $\overline{C} \wedge C_2$) participate in a triangle in $T_A$ (resp. $T_B$). This determines the edges in $(V_1,V_3)$ participating in triangles from $T_A \cup T_B$.

Set $C = C \vee (C_1 \wedge \overline{C}) \vee (C_2 \wedge \overline{C})$, and update $\overline{C}$ and the edges in $(V_1,V_3)$ accordingly. The only remaining edges in $(V_1,V_3)$ that could be involved in a triangle are those corresponding to $1$-entries in ${\overline{C}}_F$. We now need to determine which of these actually lie in a triangle. 

Our remaining problem is the following: we have a tripartite graph on vertex set $(V_1,V_2,V_3)$ with at most $f(\eps) n^2$ edges between $V_1$ and $V_3$, and each such edge lies in at most  $\eps n$ triangles. We wish to determine the edges in $(V_1,V_3)$ that participate in triangles. This problem is solved by the following theorem.

\begin{theorem}[Reporting Edges in Triangles]\label{sparseedgelisting} Let $G$ be a tripartite graph on vertex set $(V_1,V_2,V_3)$ such that there are at most $\delta n^2$ edges in $(V_1,V_3)$, and every edge of $(V_1,V_3)$ is in at most $t$ triangles. Then the set of edges in $(V_1,V_3)$ that participate in triangles can be computed in $O(\delta n^3 \log(1/\delta)/(w \log n) + n^2 t)$ time. \end{theorem}

Setting $\delta = f(\eps)$ and $t = \eps n$, \expref{Theorem}{sparseedgelisting} implies the desired time bound in \expref{Theorem}{boolmmtriangles}. The idea of the proof of \expref{Theorem}{sparseedgelisting} is to work with a new tripartite graph where the vertices have asymptotically smaller degrees, at the cost of adding slightly more nodes. This is achieved by having some nodes in our new graph correspond to \emph{small subsets of nodes} in the original tripartite graph.

\begin{proof}[Proof of~\expref{Theorem}{sparseedgelisting}] We first describe how to do the computation on a pointer machine with $w=\log n$, then describe how to modify it to work for the word RAM.


\paragraph{Graph Construction} We start by defining a new tripartite graph $G'$ on vertex set $(V_1,V_2',V_3')$.
Let $\gamma < 1/2$. $V_2'$ is obtained by partitioning the nodes of $V_2$ into $n/(\gamma \log n)$ groups of size $\gamma \log n$ each. For each group, we replace it by $2^{\gamma \log n} = n^\gamma$ nodes, one corresponding to each subset of nodes in that group. Thus $V_2'$ has $n^{1+\gamma}/(\gamma \log n)$ nodes.

$V_3'$ is also constructed out of subsets of nodes. We form $n/\ell$ groups each consisting of $\ell$ nodes in $V_3$, where $\ell = \gamma (\log n) /(\delta \log(e/\delta))$. For each group, we replace it by $\binom{\ell}{k} \leq (el/k)^k \leq n^\gamma$ nodes, one corresponding to each subset of size up to $\kappa = \gamma (\log n)/(\log (e/\delta))$. So $V_3'$ has $O(n^{1+\gamma}/\ell)$ nodes.

\smallskip
\noindent
\emph{Edges in $(V_2',V_3')$:} Put an edge between $u$ in $V_2'$ and $x$ in $V_3'$ if there is an edge $(i,j)$ in $(V_2,V_3)$ such that $i$ lies in the set corresponding to $u$, and $j$ lies in the set corresponding to $x$. For each such edge $(u,x)$, we make a list of all edges $(i,j) \in (V_2,V_3)$ corresponding to it. Observe the list for a single edge has size at most $ \gamma \log n \cdot \ell = O(\log^2 n)$. 

\smallskip
\noindent
\emph{Edges in $(V_1,V_2')$:} The edges from $v \in V_1$ to $V_2'$ are defined as follows. For each group in $V_2$ consider the neighbors of $v$ in that group. Put an edge from $v$ to the node in $V'_2$ corresponding to this subset. Each $v$ has at most $n/(\gamma \log n)$ edges to nodes in $V_2'$.

\smallskip
\noindent
\emph{Edges in $(V_1,V_3')$:} Let $v \in V_1$. For each group $g$ of $\ell$ nodes in $V_3$, let $N_{v,g}$ be the set of neighbors of $v$ in $g$. Let $d_{v,g} = |N_{v,g}|$. Partition $N_{v,g}$ arbitrarily into $t=\lceil d_{v,g}/\kappa \rceil$ subsets $s_1,\ldots,s_t$ each of size at most $\kappa$. Put edges from $v$ to $s_1,\ldots,s_t$ in $V_3'$. The number of these edges from $v$ is at most $\sum_g \lceil d_{v,g}/\kappa \rceil \leq n/\ell + d_v /\kappa$, where $d_v$ is the number of edges from $v$ to $V_3$. Since $\sum_v d_v \leq \delta n^2$, the total number of edges from $V_1$ to $V_3'$ is $O(\delta \log (1/\delta) n^2 /(\gamma \log n))$.

\paragraph{Final Algorithm} For each vertex $v \in V_1$, iterate over each pair of $v$'s neighbors $u \in V_2'$ and $x \in V_3'$. If $(u,x)$ is an edge in $G'$, output the list of edges  $(i,j)$ in $(V_2,V_3)$ corresponding to $(u,x)$, otherwise continue to the next pair. From these outputs we can easily determine the edges $(v,j)$ in $(V_1,V_3)$
that are in triangles: $(v,j)$ is in a triangle if and only if node $j$ in $V_3$ is output as an end point of some edge $(i,j) \in (V_2,V_3)$ during the loop for $v$ in $V_1$.

\smallskip
\noindent
\emph{Running Time:} The graph construction takes at most $O(n^{2+2\gamma})$. In the final algorithm, the total number of pairs $(u,w)$ in $(V_2',V_3')$ that are examined is at most \[(n/\log n) \cdot O(\delta n^2 (\log 1/\delta)/ \log n)) \leq O(\delta \log (1/\delta) n^3/\log^2 n)\,.\]

We claim that the time used to output the lists of edges is at most $O(n^2 t)$ time. A node $j$ from $V_3$ is on an output list during the loop for $v$ in $V_1$ if and only if $(v,j)$ is an edge in a triangle, with some node in $V_2$ that has a 1 in the node $i$ in $V_2'$. Since each edge from $(V_1,V_3)$ in a triangle is guaranteed to have at most $t$ witnesses in $V_2$, the node $j$ is output at most $t$ times over the loop for $v$ in $V_1$. Hence the length of all lists output during the loop for $v$ is at most $nt$, and the total time for output is at most $O(n^2 t)$.

\smallskip
\noindent
\emph{Modification for $w$-word RAM:} We now show how to replace a $\log$-speedup by a $w$-speedup with wordsize $w$.
We form $V_3'$ as above and consider the graph on $(V_1,V_2,V_3')$ (note the $V_2$ instead of $V_2'$ previously).
Recall that a node $x \in V_3'$ corresponds to a subset $V_x \subset V_3$ of at most $\kappa$ vertices (from some group of size $\ell$).
An edge $(v,x) \in (V_1,V_3')$ implies that $v$ is adjacent to every vertex in $V_x$.
For each $v \in V_1$, let $S_v$ be the $n$-bit indicator vector corresponding to neighbors of $v$ in $V_2$. For each
$x \in V_3'$, let $T_x$ denote the $n$-bit indicator vector corresponding to the union of neighbors of vertices $V_x$ in $V_2$, \ie, $T_x[i]=1$  iff some $v \in V_x$ is adjacent to $i\in V_2$.
The vectors $S_v$ and $T_x$ are stored as $n/w$ words.  With each bit $T_x[i]$ of $T_x$ such that $T_x[i]=1$, we store a pointer to a list of vertices $v \in V_x$  that are adjacent to  $i \in V_2$. It is easily checked that the overall space usage is bounded by  $O(n \cdot |V_3'| \cdot \kappa) = \tilde{O}(n^{2+\gamma})$.

The algorithms works as follows: Now for every $v$ in $V_1$ and every neighbor $x \in V'_3$ of $v$, for $i=1,\ldots,n/w$, we look up
the $i$th word $q$ from $S_v$ and the $i$th word $q'$ in the $T_x$, and compute $q \wedge q'$. If this is nonzero, then each bit location $b$ where $q\wedge q'$ has a 1 means that the node corresponding to $b$ forms a triangle with $v$ and some vertex in $V_x$. For each such bit (say at location $i$), the list associated to it consists of precisely the nodes in $V_x$ forming a triangle with $i \in V_2$ and $v$.

The running time follows as there are $O(\delta n^2 (\log 1/\delta)/ \log n))$ edges in $(V_1,V_3')$, and hence there we look up at most $O(n (\delta n^2 (\log 1/\delta)/ \log n)  (n/w)$ words $q,q'$. Since the total number of triangles is $n^2 t$, the total time spent in processing whenever
 $q \wedge q'=1$ is $O(n^2 t)$.
\end{proof}


\begin{remark} Note that we only use randomness in the BMM algorithm  to determine the pairs $(i,j)$ that have many  witnesses.  Moreover, by choosing a larger sample $R$ in the random sampling step (notice we have a lot of slack in the running time of the random sampling step), the probability of failure can be made exponentially small.\end{remark}
\clearpage
\begingroup\raggedright
%
%
%
%
%
\providecommand{\bibhead}[1]{}
%
\expandafter\ifx\csname pdfbookmark\endcsname\relax%
  \providecommand{\tocrefpdfbookmark}{}
\else\providecommand{\tocrefpdfbookmark}{%
   \hypertarget{tocreferences}{}%
   \pdfbookmark[1]{References}{tocreferences}}%
\fi

\tocrefpdfbookmark
\begin{thebibliography}{10}

\bibitem{aingworth}\bibhead{aingworth}
{\sc Donald Aingworth, Chandra Chekuri, Piotr Indyk, and Rajeev Motwani}: Fast
  estimation of diameter and shortest paths (without matrix multiplication).
\newblock {\em SIAM J. Comput.}, 28(4):1167--1181, 1999.
\newblock Preliminary version in SODA'96.
\newblock [\epfmtdoi{10.1137/S0097539796303421}]

\bibitem{bard2}\bibhead{bard2}
{\sc Martin Albrecht, Gregory Bard, and William Hart}: Algorithm 898:
  {Efficient} multiplication of dense matrices over {GF(2)}.
\newblock {\em ACM Trans. Math. Softw.}, 37(1), 2010.
\newblock [\epfmtdoi{10.1145/1644001.1644010}]

\bibitem{ADLRY}\bibhead{ADLRY}
{\sc Noga Alon, Richard~A. Duke, Hanno Lefmann, Vojtech R{\"o}dl, and Raphael
  Yuster}: The algorithmic aspects of the regularity lemma.
\newblock {\em J. Algorithms}, 16(1):80--109, 1994.
\newblock Preliminary version in FOCS'92.
\newblock [\epfmtdoi{10.1006/jagm.1994.1005}]

\bibitem{AFKS}\bibhead{AFKS}
{\sc Noga Alon, Eldar Fischer, Michael Krivelevich, and Mario Szegedy}:
  Efficient testing of large graphs.
\newblock {\em Combinatorica}, 20(4):451--476, 2000.
\newblock Preliminary version in FOCS'99.
\newblock [\epfmtdoi{10.1007/s004930070001}]

\bibitem{AFNS}\bibhead{AFNS}
{\sc Noga Alon, Eldar Fischer, Ilan Newman, and Asaf Shapira}: A combinatorial
  characterization of the testable graph properties: {It's} all about
  regularity.
\newblock {\em SIAM J. Comput.}, 39(1):143--167, 2009.
\newblock Preliminary version in STOC'06.
\newblock [\epfmtdoi{10.1137/060667177}]

\bibitem{AN04}\bibhead{AN04}
{\sc Noga Alon and Assaf Naor}: Approximating the cut-norm via {G}rothendieck's
  inequality.
\newblock {\em SIAM J. Comput.}, 35(4):787--803, 2006.
\newblock Preliminary version in STOC'04.
\newblock [\epfmtdoi{10.1137/S0097539704441629}]

\bibitem{Angluin}\bibhead{Angluin}
{\sc Dana Angluin}: The four {R}ussians' algorithm for {B}oolean matrix
  multiplication is optimal for its class.
\newblock {\em SIGACT News}, 1:29--33, 1976.
\newblock [\epfmtdoi{10.1145/1008591.1008593}]

\bibitem{ADKF70}\bibhead{ADKF70}
{\sc V.~Z. Arlazarov, E.~A. Dinic, M.~A. Kronrod, and I.~A. Faradzhev}: On
  economical construction of the transitive closure of a directed graph.
\newblock {\em Soviet Mathematics Doklady}, 11(5):1209--1210, 1970.

\bibitem{AS88}\bibhead{AS88}
{\sc Michael~D. Atkinson and Nicola Santoro}: A practical algorithm for
  {B}oolean matrix multiplication.
\newblock {\em Inf. Process. Lett.}, 29:37--38, 1988.
\newblock [\epfmtdoi{10.1016/0020-0190(88)90130-5}]

\bibitem{BKM95}\bibhead{BKM95}
{\sc Julien Basch, Sanjeev Khanna, and Rajeev Motwani}: On diameter
  verification and {B}oolean matrix multiplication, 1995.
\newblock Technical Report No. STAN-CS-95-1544, Department of Computer Science,
  Stanford University.

\bibitem{Beh}\bibhead{Beh}
{\sc F.~A. Behrend}: On sets of integers which contain no three terms in
  arithmetic progression.
\newblock {\em Proc. Nat. Acad. Sci.}, 32(12):331--332, 1946.

\bibitem{BCSX}\bibhead{BCSX}
{\sc Arnab Bhattacharyya, Victor Chen, Madhu Sudan, and Ning Xie}: Testing
  linear-invariant non-linear properties.
\newblock {\em \href{http://www.theoryofcomputing.org/articles/main}{Theory of
  Computing}}, 7:75--99, 2011.
\newblock Preliminary version in STACS'06.
\newblock [\epfmtdoi{10.4086/toc.2011.v007a006}]

\bibitem{ICALP08}\bibhead{ICALP08}
{\sc Guy~E. Blelloch, Virginia Vassilevska, and Ryan Williams}: A new
  combinatorial approach for sparse graph problems.
\newblock In {\em Proc. 35th Internat. Colloq. on Automata, Languages and
  Programming (ICALP'08)}, pp. 108--120, 2008.
\newblock [\epfmtdoi{10.1007/978-3-540-70575-8\_10}]

\bibitem{Avrim}\bibhead{Avrim}
{\sc Avrim Blum}: 2009.
\newblock Personal communication.

\bibitem{BCL}\bibhead{BCL}
{\sc Christian Borgs, Jennifer~T. Chayes, L{\'a}szl{\'o} Lov{\'a}sz, Vera~T.
  S{\'o}s, Bal{\'a}zs Szegedy, and Katalin Vesztergombi}: Graph limits and
  parameter testing.
\newblock In {\em Proc. 38th STOC}, pp. 261--270. ACM Press, 2006.
\newblock [\epfmtdoi{10.1145/1132516.1132556}]

\bibitem{Chan07}\bibhead{Chan07}
{\sc Timothy~M. Chan}: More algorithms for all-pairs shortest paths in weighted
  graphs.
\newblock In {\em Proc. 39th STOC}, pp. 590--598. ACM Press, 2007.
\newblock [\epfmtdoi{10.1145/1250790.1250877}]

\bibitem{mmexperiments}\bibhead{mmexperiments}
{\sc Siddhartha Chatterjee, Alvin~R. Lebeck, Praveen~K. Patnala, and Mithuna
  Thottethodi}: Recursive array layouts and fast matrix multiplication.
\newblock {\em IEEE Trans. on Parallel and Distributed Systems}, 13:1105--1123,
  2002.
\newblock [\epfmtdoi{10.1109/TPDS.2002.1058095}]

\bibitem{CohnUmans2}\bibhead{CohnUmans2}
{\sc Henry Cohn, Robert~D. Kleinberg, Bal{\'a}zs Szegedy, and Christopher
  Umans}: Group-theoretic algorithms for matrix multiplication.
\newblock In {\em Proc. 46th FOCS}, pp. 379--388. IEEE Comp. Soc. Press, 2005.
\newblock [\epfmtdoi{10.1109/SFCS.2005.39}]

\bibitem{CohnUmans}\bibhead{CohnUmans}
{\sc Henry Cohn and Christopher Umans}: A group-theoretic approach to fast
  matrix multiplication.
\newblock In {\em Proc. 44th FOCS}, pp. 438--449. IEEE Comp. Soc. Press, 2003.
\newblock [\epfmtdoi{10.1109/SFCS.2003.1238217}]

\bibitem{CCF}\bibhead{CCF}
{\sc Amin Coja-Oghlan, Colin Cooper, and Alan~M. Frieze}: An efficient sparse
  regularity concept.
\newblock {\em SIAM J. Discrete Math.}, 23(4):2000--2034, 2010.
\newblock [\epfmtdoi{10.1137/080730160}]

\bibitem{dor}\bibhead{dor}
{\sc Dorit Dor, Shay Halperin, and Uri Zwick}: All-pairs almost shortest paths.
\newblock {\em SIAM J. Comput.}, 29(5):1740--1759, 2000.
\newblock Preliminary version in FOCS'96.
\newblock [\epfmtdoi{10.1137/S0097539797327908}]

\bibitem{duke}\bibhead{duke}
{\sc Richard~A. Duke, Hanno Lefmann, and Vojtech R{\"o}dl}: A fast
  approximation algorithm for computing the frequencies of subgraphs in a given
  graph.
\newblock {\em SIAM J. Comput.}, 24(3):598--620, 1995.
\newblock [\epfmtdoi{10.1137/S0097539793247634}]

\bibitem{Elkin}\bibhead{Elkin}
{\sc Michael Elkin}: An improved construction of progression-free sets.
\newblock In {\em Proc. 21st Ann. ACM-SIAM Symp. on Discrete Algorithms
  (SODA'10)}, pp. 886--905, 2010.

\bibitem{federmotwani}\bibhead{federmotwani}
{\sc Tom{\'a}s Feder and Rajeev Motwani}: Clique partitions, graph compression
  and speeding-up algorithms.
\newblock {\em J. Comput. System Sci.}, 51(2):261--272, 1995.
\newblock Preliminary version in STOC'91.
\newblock [\epfmtdoi{10.1006/jcss.1995.1065}]

\bibitem{FischerMeyer}\bibhead{FischerMeyer}
{\sc Michael~J. Fischer and Albert~R. Meyer}: {B}oolean matrix multiplication
  and transitive closure.
\newblock In {\em Proc. 12th FOCS (SWAT'71)}, pp. 129--131. IEEE Comp. Soc.
  Press, 1971.
\newblock [\epfmtdoi{10.1109/SWAT.1971.4}]

\bibitem{Fox}\bibhead{Fox}
{\sc Jacob Fox}: A new proof of the graph removal lemma.
\newblock {\em Ann. of Math.}, 174(1):561--579, 2011.
\newblock [\epfmtdoi{10.4007/annals.2011.174.1.17}]

\bibitem{FKsimple}\bibhead{FKsimple}
{\sc Alan Frieze and Ravi Kannan}: A simple algorithm for constructing
  {S}zemer{\'e}di's regularity partition.
\newblock {\em Electr. J. Comb.}, 6, 1999.

\bibitem{FKfocs}\bibhead{FKfocs}
{\sc Alan~M. Frieze and Ravi Kannan}: The {R}egularity {L}emma and
  approximation schemes for dense problems.
\newblock In {\em Proc. 37th FOCS}, pp. 12--20. IEEE Comp. Soc. Press, 1996.
\newblock [\epfmtdoi{10.1109/SFCS.1996.548459}]

\bibitem{FK}\bibhead{FK}
{\sc Alan~M. Frieze and Ravi Kannan}: Quick approximation to matrices and
  applications.
\newblock {\em Combinatorica}, 19(2):175--220, 1999.
\newblock [\epfmtdoi{10.1007/s004930050052}]

\bibitem{GO95}\bibhead{GO95}
{\sc Anka Gajentaan and Mark~H. Overmars}: On a class of $o(n^2)$ problems in
  computational geometry.
\newblock {\em Computational Geometry}, 5:165--185, 1995.
\newblock [\epfmtdoi{10.1016/0925-7721(95)00022-2}]

\bibitem{GalilMargalit}\bibhead{GalilMargalit}
{\sc Zvi Galil and Oded Margalit}: All pairs shortest distances for graphs with
  small integer length edges.
\newblock {\em Inform. and Comput.}, 134:103--139, 1997.
\newblock [\epfmtdoi{10.1006/inco.1997.2620}]

\bibitem{gowers}\bibhead{gowers}
{\sc W.~T. Gowers}: Lower bounds of tower type for {S}zemer\'{e}di's uniformity
  lemma.
\newblock {\em Geom. and Funct. Anal.}, 7:322--337, 1997.
\newblock [\epfmtdoi{10.1007/PL00001621}]

\bibitem{Green}\bibhead{Green}
{\sc Ben Green}: A {S}zemer{\'e}di-type regularity lemma in abelian groups.
\newblock {\em Geom. and Funct. Anal.}, 15:340--376, 2005.
\newblock [\epfmtdoi{10.1007/s00039-005-0509-8}]

\bibitem{HMT}\bibhead{HMT}
{\sc Andr{\'a}s Hajnal, Wolfgang Maass, and Gy{\"o}rgy Tur{\'a}n}: On the
  communication complexity of graph properties.
\newblock In {\em Proc. 20th STOC}, pp. 186--191. ACM Press, 1988.
\newblock [\epfmtdoi{10.1145/62212.62228}]

\bibitem{Clique1}\bibhead{Clique1}
{\sc Alon Itai and Michael Rodeh}: Finding a minimum circuit in a graph.
\newblock {\em SIAM J. Comput.}, 7(4):413--423, 1978.
\newblock [\epfmtdoi{10.1137/0207033}]

\bibitem{kohayakawa2}\bibhead{kohayakawa2}
{\sc Y.~Kohayakawa}: {Szemer\'{e}di's} regularity lemma for sparse graphs.
\newblock In {\sc F.~Cucker and M.~Shub}, editors, {\em Foundations of
  Computational Mathematics}, pp. 216--230. Springer, 1997.

\bibitem{kohayakawa}\bibhead{kohayakawa}
{\sc Yoshiharu Kohayakawa, Vojtech R{\"o}dl, and Lubos Thoma}: An optimal
  algorithm for checking regularity.
\newblock {\em SIAM J. Comput.}, 32(5):1210--1235, 2003.
\newblock [\epfmtdoi{10.1137/S0097539702408223}]

\bibitem{Komlos}\bibhead{Komlos}
{\sc J\'{a}nos Koml\'{o}s and Mikl\'{o}s Simonovits}: {S}zemer\'{e}di's
  {R}egularity {L}emma and its applications in graph theory.
\newblock In {\em Combinatorics, Paul Erd{\H o}s is Eighty, (D. Mikl{\'o}s et.
  al, eds.), Bolyai Society Mathematical Studies}, volume~2, pp. 295--352,
  1996.

\bibitem{KSV}\bibhead{KSV}
{\sc Daniel Kr\'{a}l, Oriol Serra, and Llu\'{i}s Vena}: A removal lemma for
  systems of linear equations over finite fields.
\newblock {\em Israel J. Math.}, 187(1):193--207, 2012.
\newblock [\epfmtdoi{10.1007/s11856-011-0080-y}]

\bibitem{Lee}\bibhead{Lee}
{\sc Lillian Lee}: Fast context-free grammar parsing requires fast {B}oolean
  matrix multiplication.
\newblock {\em J. ACM}, 49:1--15, 2002.
\newblock [\epfmtdoi{10.1145/505241.505242}]

\bibitem{lingas}\bibhead{lingas}
{\sc Andrzej Lingas}: A geometric approach to {B}oolean matrix multiplication.
\newblock In {\em Proc. 13th Int. Symp. on Algorithms and Computation
  (ISAAC'02)}, pp. 501--510, 2002.
\newblock [\epfmtdoi{10.1007/3-540-36136-7\_44}]

\bibitem{LS}\bibhead{LS}
{\sc L\'{a}szl\'{o} Lov{\'a}sz and Bal\'{a}zs Szegedy}: {S}zemer\'{e}di's
  theorem for the analyst.
\newblock {\em Geom. and Funct. Anal.}, 17(1):252--270, 2007.
\newblock [\epfmtdoi{10.1007/s00039-007-0599-6}]

\bibitem{MM}\bibhead{MM}
{\sc J.~W. Moon and L.~Moser}: A matrix reduction problem.
\newblock {\em Mathematics of Computation}, 20:328--330, 1966.
\newblock [\epfmtdoi{10.1090/S0025-5718-66-99935-2}]

\bibitem{oneil}\bibhead{oneil}
{\sc Patrick~E. O'Neil and Elizabeth~J. O'Neil}: A fast expected time algorithm
  for {B}oolean matrix multiplication and transitive closure matrices.
\newblock {\em Inf. Control}, 22:132--138, 1973.
\newblock [\epfmtdoi{10.1016/S0019-9958(73)90228-3}]

\bibitem{Pan}\bibhead{Pan}
{\sc Victor~Y. Pan}: {\em How to Multiply Matrices Faster}.
\newblock Volume 179 of {\em Lecture Notes in Computer Science}.
\newblock Springer, 1984.

\bibitem{Patrascu10}\bibhead{Patrascu10}
{\sc Mihai Patrascu}: Towards polynomial lower bounds for dynamic problems.
\newblock In {\em Proc. 42nd STOC}, pp. 603--610. ACM Press, 2010.
\newblock [\epfmtdoi{10.1145/1806689.1806772}]

\bibitem{rodittyzwick}\bibhead{rodittyzwick}
{\sc Liam Roditty and Uri Zwick}: On dynamic shortest paths problems.
\newblock {\em Algorithmica}, 61(2):389--401, 2010.
\newblock Preliminary version in 12th Europ. Symp. Algor. (ESA'04).
\newblock [\epfmtdoi{10.1007/s00453-010-9401-5}]

\bibitem{RS}\bibhead{RS}
{\sc Vojt{\v e}ch R{\"o}dl and Mathias Schacht}: Property testing in
  hypergraphs and the removal lemma.
\newblock In {\em Proc. 39th STOC}, pp. 488--495. ACM Press, 2007.
\newblock [\epfmtdoi{10.1145/1250790.1250862}]

\bibitem{RuzsaSz}\bibhead{RuzsaSz}
{\sc I.~Z. Ruzsa and E.~Szemer\'{e}di}: Triple systems with no six points
  carrying three triangles.
\newblock {\em Colloquia Mathematica Societatis J\'{a}nos Bolyai}, 18:939--945,
  1978.

\bibitem{R85}\bibhead{R85}
{\sc Wojciech Rytter}: Fast recognition of pushdown automaton and context-free
  languages.
\newblock {\em Inf. Control}, 67:12--22, 1985.
\newblock Preliminary version in MFCS'84.
\newblock [\epfmtdoi{10.1016/S0019-9958(85)80024-3}]

\bibitem{S74}\bibhead{S74}
{\sc John~E. Savage}: An algorithm for the computation of linear forms.
\newblock {\em SIAM J. Comput.}, 3:150--158, 1974.
\newblock [\epfmtdoi{10.1137/0203011}]

\bibitem{schnorrsub}\bibhead{schnorrsub}
{\sc C.-P. Schnorr and C.~R. Subramanian}: Almost optimal (on the average)
  combinatorial algorithms for {B}oolean matrix product witnesses, computing
  the diameter.
\newblock In {\em Proc. 2nd Intern. Workshop Random. and Approx. Tech. Comput.
  Sci. (RANDOM'98)}, pp. 218--231, 1998.
\newblock [\epfmtdoi{10.1007/3-540-49543-6\_18}]

\bibitem{Seidel}\bibhead{Seidel}
{\sc Raimund Seidel}: On the all-pairs-shortest-path problem in unweighted
  undirected graphs.
\newblock {\em J. Comput. System Sci.}, 51(3):400--403, 1995.
\newblock [\epfmtdoi{10.1006/jcss.1995.1078}]

\bibitem{shapira}\bibhead{shapira}
{\sc Asaf Shapira}: {G}reen's conjecture and testing linear-invariant
  properties.
\newblock In {\em Proc. 41st STOC}, pp. 159--166. ACM Press, 2009.
\newblock [\epfmtdoi{10.1145/1536414.1536438}]

\bibitem{ShoshanZwick}\bibhead{ShoshanZwick}
{\sc Avi Shoshan and Uri Zwick}: All pairs shortest paths in undirected graphs
  with integer weights.
\newblock In {\em Proc. 40th FOCS}, pp. 605--614. IEEE Comp. Soc. Press, 1999.
\newblock [\epfmtdoi{10.1109/SFFCS.1999.814635}]

\bibitem{strassen}\bibhead{strassen}
{\sc Volker Strassen}: {G}aussian elimination is not optimal.
\newblock {\em Numerische Mathematik}, 13(4):354--356, 1969.
\newblock [\epfmtdoi{10.1007/BF02165411}]

\bibitem{sz}\bibhead{sz}
{\sc Endre Szemer\'edi}: Regular partitions of graphs.
\newblock In {\em Proc. Colloque Inter. CNRS (J.~C.~Bermond, J.~C.~Fournier,
  M.~Las~Vergnas and D.~Sotteau, eds.)}, pp. 399--401, 1978.

\bibitem{trevisan}\bibhead{trevisan}
{\sc Luca Trevisan}: Additive combinatorics and theoretical computer science.
\newblock In {\em SIGACT News Complexity Column 63}, 2009.

\bibitem{Valiant}\bibhead{Valiant}
{\sc Leslie~G. Valiant}: General context-free recognition in less than cubic
  time.
\newblock {\em J. Comput. System Sci.}, 10(2):308--315, 1975.
\newblock [\epfmtdoi{10.1016/S0022-0000(75)80046-8}]

\bibitem{vw12}\bibhead{vw12}
{\sc Virginia Vassilevska~Williams}: Multiplying matrices faster than
  {C}oppersmith-{W}inograd.
\newblock In {\em Proc. 44th STOC}, 2012.
\newblock To appear.

\bibitem{VWW}\bibhead{VWW}
{\sc Virginia Vassilevska~Williams and Ryan Williams}: Subcubic equivalences
  between path, matrix, and triangle problems.
\newblock In {\em Proc. 51st FOCS}, pp. 645--654. IEEE Comp. Soc. Press, 2010.
\newblock [\epfmtdoi{10.1109/FOCS.2010.67}]

\bibitem{SODA07}\bibhead{SODA07}
{\sc Ryan Williams}: Matrix-vector multiplication in sub-quadratic time (some
  preprocessing required).
\newblock In {\em Proc. 18th Ann. ACM-SIAM Symp. on Discrete Algorithms
  (SODA'07)}, pp. 995--1001. ACM Press, 2007.

\end{thebibliography}
\endgroup
\end{document}
